\chapter{Validation et déploiement local}
\label{ch:val}

\section{Vérifications automatisées}

La commande \texttt{dotnet build} sur la solution ou le projet Web doit s'exécuter sans erreur. Les migrations sont appliquées au lancement ; en cas de doute, \texttt{dotnet ef database update} peut être exécuté avec le projet d'infrastructure et le projet Web comme startup.

\section{Vérifications manuelles recommandées}

\begin{enumerate}[leftmargin=*]
  \item Lancer l'application (\texttt{dotnet run} ou \texttt{dotnet watch} depuis \texttt{GestionPortefeuille.Web}).
  \item Parcourir le CRUD actifs : création, édition, suppression, filtres.
  \item Enregistrer des transactions achat/vente : vérifier la diminution/augmentation de la liquidité et le refus des opérations invalides.
  \item Consulter le tableau de bord : cohérence des pourcentages de répartition et des gains/pertes avec le CMP.
  \item Onglet Outils : import CSV de prix, backtest sur un actif disposant d'historique.
  \item Sous Windows : si \texttt{dotnet build} échoue sur des fichiers verrouillés (\texttt{.dll}, \texttt{.pdb}), arrêter toute instance en cours (\texttt{Ctrl+C}) ou terminer le processus hôte avant de recompiler.
\end{enumerate}

\section{Limites connues}

\begin{itemize}[leftmargin=*]
  \item Historique de prix \textbf{simulé} : les conclusions du backtest et de la volatilité sont illustratives.
  \item APIs marché soumises à quotas, clés et disponibilité réseau ; le mapping crypto est limité à une liste de symboles.
  \item Pas d'authentification multi-utilisateur dans la version décrite (évolution possible).
\end{itemize}
